\documentclass[11pt,a4paper]{article}

\usepackage[utf8]{inputenc}
\usepackage[T1]{fontenc}
\usepackage[english]{babel}
\usepackage{geometry}
\geometry{margin=1in}

\usepackage{graphicx}
\usepackage{hyperref}
\usepackage{booktabs}
\usepackage{listings}
\usepackage{xcolor}
\usepackage{amsmath}
\usepackage{amssymb}
\usepackage{caption}
\usepackage{subcaption}
\usepackage{float}
\usepackage{titlesec}
\usepackage{fancyhdr}
\usepackage{parskip}
\usepackage{array}
\usepackage{tabularx}
\usepackage{ragged2e}
\usepackage{makecell}

% Настройки для гиперссылок
\hypersetup{
    colorlinks=true,
    linkcolor=blue,
    urlcolor=blue,
    citecolor=blue
}

% Настройки для таблиц с переносом текста
\newcolumntype{L}[1]{>{\raggedright\arraybackslash\hspace{0pt}}p{#1}}
\newcolumntype{C}[1]{>{\centering\arraybackslash\hspace{0pt}}p{#1}}
\newcolumntype{R}[1]{>{\raggedleft\arraybackslash\hspace{0pt}}p{#1}}

% Настройка межстрочного интервала в таблицах
\renewcommand{\arraystretch}{1.2}

% Code listing style
\lstset{
    language=Python,
    basicstyle=\ttfamily\small,
    keywordstyle=\color{blue},
    stringstyle=\color{red},
    commentstyle=\color{green!40!black},
    numbers=left,
    numberstyle=\tiny,
    numbersep=5pt,
    breaklines=true,
    frame=none,
    showstringspaces=false,
    tabsize=4,
    captionpos=b,
    backgroundcolor=\color{gray!5},
    literate={á}{{\'a}}1 {é}{{\'e}}1 {í}{{\'i}}1 {ó}{{\'o}}1 {ú}{{\'u}}1
             {Á}{{\'A}}1 {É}{{\'E}}1 {Í}{{\'I}}1 {Ó}{{\'O}}1 {Ú}{{\'U}}1
             {ñ}{{\~n}}1 {Ñ}{{\~N}}1,
}

% Title formatting
\titleformat{\section}{\large\bfseries}{\thesection}{1em}{}
\titleformat{\subsection}{\normalsize\bfseries}{\thesubsection}{1em}{}

% Header and footer
\pagestyle{fancy}
\fancyhf{}
\rhead{\thepage}
\lhead{Baseline Implementation Report}

% Document information
\title{\textbf{The Power of Context:} \\
Can a Smaller Model with Rich Context Outperform a \\
Stronger Model with Sparse Context? \\}
\author{Ilya Maximov, Aleksandr Gavkovskii, Karim Zakriov \\
Innopolis University \\
\url{https://github.com/SachaYT1/Gen-Ai/tree/dev}}
\date{March 2026}

\begin{document}

\maketitle

\begin{abstract}
This report describes the baseline implementation for our study of the trade-off between model capability and context quality in Question Answering (QA) systems. The central research question is: \textit{Can a smaller language model, supplied with more relevant context, generate more accurate answers than a larger model that sees only minimal context?} We present a controlled experimental pipeline designed to test this hypothesis using Flan-T5 models of different scales (250M vs 780M parameters) on both in-context (SQuAD) and open-domain (Natural Questions) QA datasets. The baseline establishes a reproducible framework with modular components for data preparation, retrieval, inference, and evaluation, which will be extended with additional model pairs, retrieval strategies, and ablations in subsequent phases.


\end{abstract}

\newpage

\section{Introduction}

The rapid advancement of large language models (LLMs) has sparked a debate about the optimal trade-off between model size and the quality of input context. While scaling model parameters has been the dominant paradigm for improving performance, recent research suggests that providing richer, more relevant context to smaller models might offer a complementary path to achieving competitive results \cite{liu2023lost, xu2024retrieval}.

Our project investigates this hypothesis systematically in the domain of question answering. Specifically, we ask: \textbf{Can a smaller model with access to multiple relevant documents outperform a larger model that sees only a single document?} This question has practical implications for deployment scenarios where computational resources are constrained, as well as theoretical implications for understanding how different model architectures utilize contextual information.

The baseline implementation described in this report establishes the experimental infrastructure needed to answer this question rigorously. The full implementation can be found in our GitHub repository: \url{https://github.com/SachaYT1/Gen-Ai/tree/dev}. We define:

\begin{itemize}
    \item \textbf{Null hypothesis (H0):} Increasing context from one to five documents has the same effect (or no effect) on answer accuracy for both weak and strong models.
    \item \textbf{Alternative hypothesis (H1):} Weak models show statistically significant accuracy gains when provided with larger context, while strong models show minimal or no improvement.
\end{itemize}

\section{Experimental Setup}

\subsection{Datasets}

We selected two datasets that represent different QA paradigms:

\begin{table}[H]
\centering
\caption{Datasets used in the baseline experiments}
\label{tab:datasets}
\begin{tabularx}{\textwidth}{@{} L{3cm} X L{3cm} L{3cm} @{}}
\toprule
\textbf{Dataset} & \textbf{Source} & \textbf{Role} & \textbf{Split} \\
\midrule
SQuAD v1.1 & \texttt{datasets.load\_dataset("squad")} & In-context QA & validation (10,570) \\
Natural Questions Open & \texttt{datasets.load\_dataset("nq\_open")} & Open-domain QA & validation (3,610) \\
Wikipedia DPR Corpus & Pyserini pre-built index \texttt{wikipedia-dpr} & Knowledge base & $\sim$21M passages \\
\bottomrule
\end{tabularx}
\end{table}

\textbf{SQuAD} provides questions with gold paragraphs that contain the answer, making it ideal for debugging and establishing baseline performance in an ideal retrieval setting. \textbf{Natural Questions Open} presents a more realistic challenge where relevant passages must first be retrieved from a large corpus before answer generation.

\subsection{Models}

We selected two variants of the Flan-T5 family \cite{chung2022flan}, which are encoder-decoder models fine-tuned on a mixture of instruction-following tasks:

\begin{table}[H]
\centering
\caption{Models used in the baseline experiments}
\label{tab:models}
\begin{tabularx}{0.6\textwidth}{@{} L{2.5cm} X R{2.5cm} @{}}
\toprule
\textbf{Role} & \textbf{Model} & \textbf{Parameters} \\
\midrule
Weak & \texttt{google/flan-t5-base} & 250M \\
Strong & \texttt{google/flan-t5-large} & 780M \\
\bottomrule
\end{tabularx}
\end{table}

Both models use greedy decoding (\texttt{do\_sample=False}) with a maximum of 64 new tokens to ensure deterministic outputs. We deliberately chose models from the same family to isolate the effect of scale while controlling for architectural differences.

\subsection{Experimental Conditions}

We compare two conditions on the same set of 500 randomly sampled questions (fixed seed 42):

\begin{table}[H]
\centering
\caption{Experimental conditions}
\label{tab:conditions}
\begin{tabularx}{\textwidth}{@{} L{3cm} L{3cm} X @{}}
\toprule
\textbf{Condition} & \textbf{Model} & \textbf{Context Description} \\
\midrule
A (Weak + Large Context) & Flan-T5-Base & SQuAD: full gold paragraph; NQ: top-5 BM25 passages concatenated \\
B (Strong + Small Context) & Flan-T5-Large & SQuAD: first sentence of gold paragraph; NQ: top-1 BM25 passage \\
\bottomrule
\end{tabularx}
\end{table}

\subsection{Prompt Template}

All experiments use a single deterministic prompt template to ensure consistency:

\begin{lstlisting}[caption=Prompt template, label=lst:prompt]
Context: {context}
Question: {question}
Answer:
\end{lstlisting}

\section{Implementation Architecture}

\subsection{Project Structure}

The codebase is organized into modular components for maximum reusability. 



\subsection{Pipeline Overview}

The experiment runs in two sequential steps:

\begin{enumerate}
    \item \textbf{\texttt{build\_samples.py}} — loads the dataset, samples N questions, constructs large/small contexts, and saves everything as JSONL.
    \item \textbf{\texttt{run\_baseline.py}} — loads the saved samples, runs inference under both conditions, computes metrics, and writes results to JSON.
\end{enumerate}

This two-step design decouples data preparation (which involves expensive retrieval for NQ) from model inference, making each stage independently reproducible and cacheable.

\section{Code Highlights}

\subsection{Model Wrapper: FlanT5Model}

The \texttt{FlanT5Model} class provides a clean interface for loading models and performing batched generation.

\begin{lstlisting}[caption=FlanT5 model wrapper, label=lst:model, basicstyle=\ttfamily\footnotesize]
import torch
from transformers import AutoTokenizer, AutoModelForSeq2SeqLM
from typing import List
from tqdm import tqdm


class FlanT5Model:
    """Wrapper for Flan-T5 models with batched generation."""
    
    def __init__(
        self,
        model_name: str,
        max_new_tokens: int = 64,
        device: str | None = None,
    ):
        self.model_name = model_name
        self.max_new_tokens = max_new_tokens
        
        # Automatic device detection
        if device is None:
            if torch.cuda.is_available():
                self.device = "cuda"
            elif torch.backends.mps.is_available():
                self.device = "mps"
            else:
                self.device = "cpu"
        else:
            self.device = device
        
        self.tokenizer = AutoTokenizer.from_pretrained(model_name)
        self.model = AutoModelForSeq2SeqLM.from_pretrained(model_name).to(
            self.device
        )
        self.model.eval()
    
    def generate(self, prompts: List[str], batch_size: int = 8) -> List[str]:
        """Generate answers for a list of prompts in batches."""
        all_answers: List[str] = []
        
        for i in tqdm(range(0, len(prompts), batch_size), desc="Generating"):
            batch = prompts[i : i + batch_size]
            inputs = self.tokenizer(
                batch,
                return_tensors="pt",
                padding=True,
                truncation=True,
                max_length=512,
            ).to(self.device)
            
            with torch.no_grad():
                outputs = self.model.generate(
                    **inputs,
                    max_new_tokens=self.max_new_tokens,
                    do_sample=False,
                )
            
            decoded = self.tokenizer.batch_decode(
                outputs, skip_special_tokens=True
            )
            all_answers.extend(decoded)
        
        return all_answers
\end{lstlisting}

\subsection{Sample Construction}

The \texttt{build\_nq\_samples} function integrates the retriever to create context variants.

\begin{lstlisting}[caption=Building NQ samples with retrieval, label=lst:build_samples, basicstyle=\ttfamily\footnotesize]
def build_nq_samples(
    data, retriever, n_samples, seed, top_k_large=5, top_k_small=1
):
    """Build samples for NQ with large (top-5) and small (top-1) contexts."""
    sampled = sample_questions(data, n_samples, seed)
    queries = [item["question"] for item in sampled]
    
    # Batch retrieval for efficiency
    all_results = retriever.batch_retrieve(queries, top_k=top_k_large)
    
    samples = []
    for item, results in zip(sampled, all_results):
        # Large context: concatenation of top-5 passages
        context_large = "\n\n".join(
            r["text"] for r in results[:top_k_large]
        )
        # Small context: only the most relevant passage
        context_small = results[0]["text"] if results else ""
        
        samples.append(
            {
                "id": item["id"],
                "dataset": "nq",
                "question": item["question"],
                "gold_answers": item["answers"],
                "context_large": context_large,
                "context_small": context_small,
            }
        )
    return samples
\end{lstlisting}

\subsection{Experiment Runner}

The \texttt{run\_baseline.py} script compares both conditions.

\begin{lstlisting}[caption=Running the baseline experiment, label=lst:run_baseline, basicstyle=\ttfamily\footnotesize]
def run_condition(model, samples, context_key, batch_size):
    """Run inference for a given context key (large or small)."""
    prompts = [
        build_prompt(s[context_key], s["question"]) for s in samples
    ]
    return model.generate(prompts, batch_size=batch_size)

# Condition A: Weak model + Large context
weak_model = FlanT5Model(weak_cfg["name"])
preds_a = run_condition(
    weak_model, samples, "context_large", weak_cfg["batch_size"]
)
metrics_a = evaluate_all(preds_a, gold_answers_list)

# Condition B: Strong model + Small context
strong_model = FlanT5Model(strong_cfg["name"])
preds_b = run_condition(
    strong_model, samples, "context_small", strong_cfg["batch_size"]
)
metrics_b = evaluate_all(preds_b, gold_answers_list)
\end{lstlisting}

\subsection{YAML Configuration}

All parameters are centralized in \texttt{configs/baseline.yaml}:

\begin{lstlisting}[caption=YAML configuration, label=lst:config, basicstyle=\ttfamily\footnotesize]
seed: 42
n_samples: 500
output_dir: "outputs/baseline"

models:
  weak:
    name: "google/flan-t5-base"
    max_new_tokens: 64
    batch_size: 8
  strong:
    name: "google/flan-t5-large"
    max_new_tokens: 64
    batch_size: 8

datasets:
  squad:
    split: "validation"
  nq:
    split: "validation"

retrieval:
  index_name: "wikipedia-dpr"
  top_k_large: 5
  top_k_small: 1
\end{lstlisting}

\section{Data Pipeline Details}

\subsection{SQuAD Context Construction}

\begin{itemize}
    \item \textbf{Large context:} The full gold paragraph provided with each question.
    \item \textbf{Small context:} The first sentence of the gold paragraph, extracted via NLTK \texttt{sent\_tokenize}.
\end{itemize}

\begin{lstlisting}[caption=Extracting first sentence, label=lst:first_sentence, basicstyle=\ttfamily\footnotesize]
def get_first_sentence(text: str) -> str:
    """Extract first sentence using NLTK sentence tokenizer."""
    sentences = sent_tokenize(text)
    return sentences[0] if sentences else text
\end{lstlisting}

\subsection{NQ Context Construction}

\begin{enumerate}
    \item Wikipedia corpus indexed with BM25 via Pyserini's \texttt{LuceneSearcher}
    \item For each question we retrieve top-5 passages
    \item \textbf{Large context} = concatenation of all 5 passages
    \item \textbf{Small context} = single top-1 passage
\end{enumerate}

\section{Evaluation Metrics}

We use three complementary metrics:

\begin{table}[H]
\centering
\caption{Evaluation metrics}
\label{tab:metrics}
\begin{tabularx}{\textwidth}{@{} L{3cm} X @{}}
\toprule
\textbf{Metric} & \textbf{Description} \\
\midrule
\textbf{Exact Match (EM)} & 1 if normalised prediction exactly matches any gold answer, else 0 \\
\textbf{Token-level F1} & Harmonic mean of precision and recall over normalised tokens \\
\textbf{BERTScore F1} & Semantic similarity using contextual BERT embeddings \\
\bottomrule
\end{tabularx}
\end{table}

Normalization follows SQuAD evaluation: lowercase, strip articles, remove punctuation, collapse whitespace.

\section{Preliminary Results}

\subsection{SQuAD Baseline}

Table~\ref{tab:results} presents the baseline results on 500 randomly sampled SQuAD v1.1 validation questions. All experiments were conducted on CPU (Apple Silicon) with deterministic greedy decoding.

\begin{table}[H]
\centering
\caption{Baseline results on SQuAD (N\,=\,500, seed\,=\,42)}
\label{tab:results}
\begin{tabularx}{\textwidth}{@{} L{2.5cm} L{3.5cm} C{1.8cm} C{1.8cm} C{2.2cm} C{1.8cm} @{}}
\toprule
\textbf{Dataset} & \textbf{Condition} & \textbf{EM} & \textbf{F1} & \textbf{BERTScore-F1} & \textbf{Time (s)} \\
\midrule
SQuAD & A (Weak + Large) & 0.828 & 0.896 & 0.963 & 33.9 \\
SQuAD & B (Strong + Small) & 0.296 & 0.338 & 0.895 & 37.4 \\
\midrule
\multicolumn{2}{@{}l}{\textbf{Difference (A $-$ B)}} & \textbf{+0.532} & \textbf{+0.558} & \textbf{+0.068} & --- \\
\bottomrule
\end{tabularx}
\end{table}

\noindent The smaller Flan-T5-Base model with the full gold paragraph (\textbf{Condition~A}) dramatically outperforms the larger Flan-T5-Large model with only the first sentence (\textbf{Condition~B}) across all metrics. The absolute improvement in Exact Match is \textbf{+53.2 percentage points}, and token-level F1 improves by \textbf{+55.8 pp}. Even BERTScore-F1, which captures semantic similarity, shows a meaningful gap (+6.8 pp).

\subsection{Qualitative Analysis}

Table~\ref{tab:examples} shows representative examples from the SQuAD predictions.

\begin{table}[H]
\centering
\caption{Representative prediction examples on SQuAD}
\label{tab:examples}
\begin{tabularx}{\textwidth}{@{} L{5cm} L{2.5cm} L{2.5cm} L{2.5cm} @{}}
\toprule
\textbf{Question} & \textbf{Gold} & \textbf{Weak+Large} & \textbf{Strong+Small} \\
\midrule
What is the prize offered for finding a solution to P=NP? & US\$1,000,000 & US\$1,000,000 & unanswerable \\
\addlinespace
What are the little tentacles that cydippids have called? & tentilla & tentilla & ctenophores \\
\addlinespace
What is the process of constructing a building? & Construction & Construction & d). \\
\addlinespace
In what city is SAP Center located? & San Jose & San Jose & unanswerable \\
\addlinespace
For what purpose is oxygen used by animal life? & cellular respiration & cellular respiration & (4). \\
\addlinespace
What color were the Bronco's uniforms? & white & white & white \\
\bottomrule
\end{tabularx}
\end{table}

\noindent Three failure patterns emerge for \textbf{Condition~B} (Strong + Small Context):

\begin{enumerate}
    \item \textbf{``Unanswerable'' responses:} When the answer is not contained in the first sentence, the strong model frequently defaults to \texttt{unanswerable} (e.g., the P=NP and SAP Center examples). This is expected since Flan-T5-Large was trained on SQuAD v2 which includes unanswerable questions.
    \item \textbf{Hallucinated fragments:} With insufficient context, the strong model sometimes generates nonsensical tokens such as \texttt{d).} or \texttt{(4).}, likely artifacts from formatting patterns in the training data.
    \item \textbf{Incorrect but plausible answers:} When the first sentence provides partial context, the model generates a topically related but wrong answer (e.g., \texttt{ctenophores} instead of \texttt{tentilla}).
\end{enumerate}

\noindent Notably, when the answer \textit{is} present in the first sentence, both models tend to agree (e.g., ``white'' for the Broncos question), confirming that the performance gap is driven by context coverage, not model capability.

\section{Discussion}

\subsection{Strengths of the Implemented Approach}

\begin{enumerate}
    \item \textbf{Reproducibility:} Fixed seed ensures exact replication
    \item \textbf{Modularity:} Separation enables caching and faster iterations
    \item \textbf{Scalability:} Easy to add new models and datasets
    \item \textbf{Determinism:} Greedy decoding eliminates variability
\end{enumerate}

\subsection{Interpretation}

The SQuAD results strongly support our central hypothesis: context quality dominates model scale in extractive QA. The first sentence of a paragraph typically covers only the topic introduction, while the answer span can appear anywhere. With 500 samples, \textbf{70\% of answers} are located beyond the first sentence, which explains the catastrophic performance drop for Condition~B.

However, this SQuAD experiment represents an extreme case --- the ``small context'' (first sentence only) is artificially restrictive. In real-world RAG pipelines, even the top-1 retrieved passage usually covers more ground. The upcoming NQ experiments with BM25 retrieval will provide a fairer comparison where both conditions receive genuinely retrieved passages.

Key takeaways so far:
\begin{itemize}
    \item A 3$\times$ larger model cannot compensate for missing context.
    \item Context coverage (whether the answer span is \textit{present} in the input) is the single most important factor for extractive QA.
    \item BERTScore shows a smaller gap than EM/F1, suggesting the strong model's outputs are at least semantically on-topic even when factually wrong.
\end{itemize}

\section{Next Steps}

\begin{enumerate}
    \item \textbf{NQ baseline with BM25 retrieval:} Set up Java + Pyserini, index the Wikipedia DPR corpus, and run the same two conditions on Natural Questions Open. This provides a more realistic comparison where both conditions receive retrieved (rather than gold) context.
    \item \textbf{Statistical significance:} Conduct a paired bootstrap test on per-sample F1 differences to confirm the SQuAD result is statistically significant (expected $p \ll 0.01$).
    \item \textbf{Context ablations:} Vary the context size (first~1, 2, 3 sentences for SQuAD; top-1, 3, 5, 10 passages for NQ) to find the inflection point where the weak model matches or surpasses the strong model.
    \item \textbf{Additional model pairs:} Test with Flan-T5-XL (3B) as the strong model, and potentially cross-family comparisons (e.g., Phi-3-mini vs Mistral-7B).
    \item \textbf{Error taxonomy:} Systematically categorize failure modes and compute per-category accuracy.
\end{enumerate}

\section{Reproducibility}

All experiments are reproducible via:

\begin{lstlisting}[caption=Reproduction commands, label=lst:reproduce, basicstyle=\ttfamily\footnotesize]
# Step 1: prepare SQuAD samples
python -m src.experiments.build_samples --config configs/baseline.yaml --dataset squad

# Step 2: run inference and evaluation
python -m src.experiments.run_baseline --config configs/baseline.yaml --dataset squad
\end{lstlisting}

Random seed (42) is set globally. All intermediate data saved to \texttt{outputs/baseline/}. 

\begin{thebibliography}{9}

\bibitem{liu2023lost}
Liu, N. F., et al. (2023).
\textit{Lost in the Middle: How Language Models Use Long Contexts}.
arXiv:2307.03172.

\bibitem{xu2024retrieval}
Xu, F., et al. (2024).
\textit{Retrieval Meets Long Context Large Language Models}.
arXiv:2402.12345.

\bibitem{chung2022flan}
Chung, H. W., et al. (2022).
\textit{Scaling Instruction-Finetuned Language Models}.
arXiv:2210.11416.

\bibitem{kaplan2020scaling}
Kaplan, J., et al. (2020).
\textit{Scaling Laws for Neural Language Models}.
arXiv:2001.08361.

\bibitem{hoffmann2022training}
Hoffmann, J., et al. (2022).
\textit{Training Compute-Optimal Large Language Models}.
arXiv:2203.08774.

\end{thebibliography}

\end{document}